\chapter{Literature Review}
\section{Review scope and selection}
We reviewed research in four areas: learner modelling, tutoring and hint provision for programming, LLM-based programming assistants, and multi-agent or conversational tutors. We searched the ACM Digital Library, the ACL Anthology, arXiv, the AAAI and IJCAI proceedings, and the Springer and Elsevier journal platforms. The search used combinations of the keywords \emph{intelligent tutoring system}, \emph{knowledge tracing}, \emph{programming education}, \emph{hint generation}, \emph{LLM programming assistant}, \emph{help seeking}, \emph{scaffolding} and \emph{multi-agent tutoring}.

A work was selected if it met three conditions. It had to be a primary research article in a peer-reviewed venue or a widely used preprint. It had to address learner modelling, help provision, programming practice or LLM-based tutoring. It also had to inform a specific design decision in CSMASF. The publication period is 1995 to 2025. We reviewed thirteen works. Nine were listed in the initial project submission, and we added four to cover the foundations of knowledge tracing \cite{corbett1995}, help-seeking research \cite{aleven2016}, guardrails for LLM help \cite{liffiton2023} and multi-agent LLM tutoring \cite{schmucker2024}.

We checked the bibliographic details of each work against its publisher or index record (DOI, ACL Anthology, IJCAI, arXiv, Crossref or OpenAlex). Findings are reported as stated in each work's abstract or published record. Numerical values are given only where they appear there or, for Deep Knowledge Tracing, in its results table. Where a detail is not stated in the material we consulted, we say so and do not assume it.

\section{Descriptive review of literature works}
The works are presented in order of publication year.

\clearpage
\subsection{Knowledge tracing}
Corbett and Anderson's ``Knowledge tracing: Modeling the acquisition of procedural knowledge'' (1995) \cite{corbett1995} asks how a tutor can model a student's knowledge state as it changes during skill acquisition. The setting is the ACT Programming Tutor, with which students learn to write short programs. The tutor has a cognitive model of programming made of production rules (the ``ideal student model'', grounded in ACT-R theory). In \emph{knowledge tracing}, while the student works, the tutor maintains for every rule an estimate of the probability that the student has learned it. Each rule is assumed to be in one of two states, learned or not learned, and the model has parameters for initial knowledge, the chance of learning at each opportunity, guessing and slipping. The estimates drive mastery learning, in which the tutor presents an individualised sequence of exercises until each rule is mastered. To assess knowledge tracing, the authors used its estimates to predict students' performance on test exercises, and they report that the model predicts that performance well. The approach needs a hand-authored rule model for each domain, and that is costly to build. We took the mastery model of our Cognitive State agent directly from this paper. The agent uses the same two-state update with fixed guess, slip and learning parameters, because it is interpretable and needs no training data.

\subsection{Deep knowledge tracing}
``Deep Knowledge Tracing'' by Piech et al. (2015) \cite{piech2015} takes up the same task, modelling a student's knowledge as they work through exercises, for large online learning platforms. The authors use recurrent neural networks, specifically long short-term memory (LSTM) networks. These learn from response sequences and need no expert annotation of the skills each exercise requires. The study uses simulated students, Khan Academy mathematics data (47,495 students, 69 exercise tags, about 1.4 million answers) and the ASSISTments benchmark (15,931 students, 124 exercise tags, about 526 thousand answers). The model is evaluated on predicting the next response, measured by the area under the ROC curve (AUC) with five-fold cross-validation on the real datasets. On the Khan data the reported AUC is 0.85, against 0.68 for standard Bayesian Knowledge Tracing (BKT). On ASSISTments it is 0.86 against the best previously reported 0.69, a 25\% gain, and on simulated data the model reached 0.82 against 0.54 for BKT. For tutoring, the main limitation is that the model predicts performance but does not itself decide what help to give. Its internal state is also less directly interpretable than BKT's parameters. BKT, which is the baseline in this paper, is the model we kept, because it is interpretable. A learned model becomes an option for us once real interaction data exists.

\clearpage
\subsection{DeepTutor}
Rus, Niraula and Banjade (2015) describe their system in ``DeepTutor: An Effective, Online Intelligent Tutoring System That Promotes Deep Learning'' \cite{rus2015}. Their question is how educational technology can deliver instruction tailored to each learner. DeepTutor is a conversational intelligent tutoring system delivered as a web application, so learners can use it at any time from a browser-connected device. It was used in large-scale experiments with high-school students, and the authors report these experiments as evidence that conversational intelligent tutoring can work at scale over the web. The published abstract does not give detailed metrics. For our purposes the system has two limitations: it predates current LLMs, and it is not designed for programming practice with code execution. What we took from it is that dialogue-based tutoring can be delivered at scale through the web. This supports CodePath's browser-based design and its conversational reflection step.

\subsection{Help seeking with intelligent tutors}
In ``Help Helps, But Only So Much: Research on Help Seeking with Intelligent Tutoring Systems'' \cite{aleven2016}, Aleven, Roll, McLaren and Koedinger (2016) study how learners seek and use on-demand help in intelligent tutoring systems. They built the Help Tutor, an agent that gave students real-time, in-context feedback on their help-seeking behaviour. Behind it is a knowledge-engineered, rule-based executable model of help seeking. The Help Tutor was evaluated in a classroom study of students' help-seeking behaviour and domain learning outcomes, and the paper adds a theoretical analysis and a review of later empirical work. The feedback did make students use on-demand help more deliberately, and this effect persisted after the feedback stopped. It did not improve domain-level learning. The authors now place less weight on on-demand, principle-based help during problem solving, although they still recommend that tutors include it. They describe helping students become better self-regulated learners as an open challenge. Three decisions in our project come from this paper. Help that the learner did not request is offered and is shown only if the learner accepts it. Help climbs only when the learner asks. And the help itself is designed to keep the learner thinking.

\clearpage
\subsection{Project CodeNet}
Puri et al. (2021) present ``Project CodeNet: A Large-Scale AI for Code Dataset for Learning a Diversity of Coding Tasks'' \cite{puri2021} in response to the lack of large, diverse datasets for training AI models on coding tasks. CodeNet is a dataset of over 14 million code samples, about 500 million lines of code, in 55 programming languages. The samples are submissions to online programming problems. They are annotated to support code similarity and classification, translation between languages, and runtime and memory improvement, and sample input and output tests are provided for 98.5\% of them. The authors supply preprocessing tools that turn source code into model inputs. They evaluate the dataset by reporting baseline results for code classification and code similarity. CodeNet contains submissions and their outcomes but no records of tutoring, hints or learner states, so by itself it does not support learner modelling or tutoring. For us it is a candidate data source for offline evaluation. CodePath includes a CodeNet importer that aggregates problem statistics, and we tested it on a fixture, but the archive itself was not imported in this phase.

\subsection{CodeHelp}
``CodeHelp: Using Large Language Models with Guardrails for Scalable Support in Programming Classes'' by Liffiton, Sheese, Savelka and Denny (2023) \cite{liffiton2023} is about large computing classes, where students need timely help but should not be encouraged to over-rely on LLMs. CodeHelp is an LLM-based tool with guardrails, so it helps students without revealing solutions. It has instructor-facing features and uses a sequence of prompting strategies to keep its outputs suitable for students. The tool was used in a first-year computer and data science course with 52 students, and their interactions were collected over 12 weeks. The evaluation covers usage patterns and student perceptions and includes the instructor's reflections. Students valued the tool's availability and its help in resolving errors. For instructors it was easy to deploy, and it complemented their support without replacing it. The published description does not include a per-learner model that adapts the amount of help over time. CodePath's rule that help never contains a solution comes from this work, and so does its validator, which checks every hint before it is shown.

\clearpage
\subsection{CodeAid}
Kazemitabaar et al. (2024) report on a classroom deployment in ``CodeAid: Evaluating a Classroom Deployment of an LLM-based Programming Assistant that Balances Student and Educator Needs'' \cite{kazemitabaar2024}. They wanted an LLM assistant that gives students helpful, technically correct support without revealing code solutions. CodeAid answers conceptual questions, produces pseudo-code with explanations and suggests fixes for incorrect student code. It was deployed for 12 weeks in a course with 700 students. The authors analysed about 8,000 uses together with weekly surveys and 22 student interviews, and then interviewed eight educators. From this they derive four design considerations: exploit what AI uniquely offers, simplify query formulation while promoting cognitive engagement, avoid direct answers while encouraging motivated learning, and keep the assistant transparent and controllable for students. The published description does not include a continuous learner model or retrieval of what helped similar learners. Two CodePath features follow from this work. The typed-question feature answers from what the selected lines actually did, and the ``Why this hint?'' explanation is there for transparency.

\subsection{PyTutor}
Yang, Lin, Lin and Ogata (2024) test whether a ChatGPT-based tutoring system can help beginners learn Python, in ``Enhancing Python Learning with PyTutor: Efficacy of a ChatGPT-Based Intelligent Tutoring System in Programming Education'' \cite{yang2024}. PyTutor gives continuous guidance, problem-solving hints and code explanations. Each question has a structured set of hints covering pseudocode, cloze, basic and advanced coding solutions. In an eleven-week experiment 35 students used PyTutor and 36 did not, and the authors compared engagement, completion and success rates in in-class and after-class exercises. PyTutor was especially effective for students with weak foundations, who showed higher engagement, completion and success. The authors also identify a risk of over-reliance that may hinder independent problem-solving skills, and they recommend balanced use. A single LLM generates the help, and the hint set includes coding solutions. We took two things from this study. It supports a graded ladder of help. Its over-reliance finding is also the reason for two CodePath rules: no level contains a solution, and unrequested help is only offered.

\clearpage
\subsection{Integrating AI tutors in a programming course}
``Integrating AI Tutors in a Programming Course'' by Ma, Krone Martins and Lopes (2024) \cite{ma2024} deals with course-specific AI tutoring for programming homework. Their system, RAGMan, is LLM-powered and supports course-specific and homework-specific tutors through retrieval-augmented generation and strict instructions. It was offered in an introductory programming course with 455 students, as five homework-specific tutors that students could use if they wished. The authors evaluate usage, answer accuracy within the intended scope and student perceptions. About half of the students used the tutors, mostly for legitimate homework questions. Within scope the tutors answered accurately 98\% of the time. Of the students who used them, 78\% said the tutors helped their learning, and students valued the non-judgmental environment. The retrieval draws on course material, and the outcomes are largely self-reported. Our retrieval design differs on the first point. RAGMan retrieves course content, whereas CodePath retrieves what helped in similar \emph{learner situations}.

\subsection{Ruffle\&Riley}
Schmucker, Xia, Azaria and Mitchell (2024) look at how LLMs can lower the cost of authoring conversational tutoring systems, in ``Ruffle\&Riley: Insights from Designing and Evaluating a Large Language Model-Based Conversational Tutoring System'' \cite{schmucker2024}. Their system automatically generates an editable tutoring script from a lesson text. It then runs a learning-by-teaching format in which two LLM agents play a student and a professor. The evaluation used biology lessons in two between-subject online studies with 200 participants in total, and compared the system with simpler question-answering chatbots and a reading activity. Users found the system engaging and helpful and felt they understood the material, and they spent more time on the activity. Their short-term learning gains, however, did not differ significantly from those of the reading activity. The domain is not programming. The paper matters to us for two reasons. It shows that several cooperating agents can structure a tutoring interaction. It also cautions that engagement and perceived helpfulness are not evidence of learning, which is why this report does not claim learning gains without a learner study.

\clearpage
\subsection{CoderAgent}
Zhan et al. (2025), in ``CoderAgent: Simulating Student Behavior for Personalized Programming Learning with Large Language Models'' \cite{zhan2025}, start from a practical difficulty. Real programming data is scarce, and offline evaluation of personalised programming tutoring often does not match real learning. They propose CoderAgent, an LLM-based agent inspired by the ACT-R cognitive architecture, which simulates students' iterative programming processes. It models both programming knowledge and coding ability, and its Programming Tree of Thought breaks each step into why, how, where and what. The study uses real-world datasets, including a Python subset of CodeNet and a college-level Java dataset (CSEDM). The simulation is evaluated on four predictions: a student's next modification intention, its location, the next submission and whether that submission will be accepted. The measures are accuracy, CodeBLEU and AUC. The authors report that CoderAgent gives interpretable insights into learning trajectories and simulates students accurately. It simulates learners and does not tutor them in real time. We see it as a possible way to test tutoring policies offline before a learner study.

\subsection{TRAVER}
Wang, Dai, Zhang, Ma, Li and Chai (2025) study how LLM tutoring agents can proactively guide learners to complete predefined coding tasks, in ``Training Turn-by-Turn Verifiers for Dialogue Tutoring Agents: The Curious Case of LLMs as Your Coding Tutors'' \cite{wang2025}. Their method, Trace-and-Verify (TRAVER), has two parts. Knowledge tracing estimates the student's knowledge state, and turn-by-turn verification keeps the tutor's guidance on track. The paper also introduces DICT, an automatic evaluation protocol based on controlled student simulation and code-generation tests. The authors report that coding tutoring is challenging and that TRAVER achieves a significantly higher success rate. The evaluation relies on simulated students. The approach includes neither retrieval of similar learners' experience nor an explicit decision about how much help to give. CodePath combines a learner-state estimate with a verification step in a similar way: a validator checks every hint before it is shown.

\clearpage
\subsection{Survey of LLM applications in programming education}
The last work is a survey. In ``A Survey of LLM-Based Applications in Programming Education: Balancing Automation and Human Oversight'' \cite{pitts2025}, Pitts, Hridi and Lekshmi Narayanan (2025) ask how LLMs can scale timely, personalised support for novice programmers when instructor time is limited. They survey recent LLM applications in three areas: formative code feedback, assessment and knowledge modelling. The published abstract does not state how many works were reviewed. The authors identify recurring design patterns. They report that interventions work best when educator expertise complements model output through human-in-the-loop oversight, scaffolding and evaluation, while fully automated approaches often fail to capture the pedagogical nuances of programming education. They call for greater transparency, closer alignment with pedagogy and systems that adapt to varied learning contexts. Being a survey, the paper does not propose or evaluate a system. Three of CodePath's design choices follow its conclusions. The policy, and not an LLM, controls the help decisions. Every decision is explained. Teachers get a dashboard.

\clearpage\begin{landscape}
\section{Comparison of reviewed literature}
Table~\ref{tab:literature} compares the reviewed works. ``Not reported'' means that the detail is not stated in the material we consulted, and ``Not applicable'' means that it does not apply to that work.

\begingroup\small\setlength{\tabcolsep}{4pt}\renewcommand{\arraystretch}{1.25}
\begin{longtable}{|R{2.2cm}|R{2.6cm}|R{3.5cm}|R{2.7cm}|R{3.0cm}|R{2.9cm}|R{3.3cm}|}
\caption{Comparison of reviewed literature}\label{tab:literature}\\\hline
\textbf{Work and citation} & \textbf{Problem or objective} & \textbf{Approach and contribution} & \textbf{Study context} & \textbf{Evaluation criteria} & \textbf{Key findings} & \textbf{Limitations and gap}\\\hline\endfirsthead
\textbf{Work and citation} & \textbf{Problem or objective} & \textbf{Approach and contribution} & \textbf{Study context} & \textbf{Evaluation criteria} & \textbf{Key findings} & \textbf{Limitations and gap}\\\hline\endhead
Work 01 \cite{corbett1995} Corbett and Anderson, 1995 & Model changing knowledge during skill learning & Knowledge tracing: probability each rule is learned, used for mastery learning & ACT Programming Tutor; students writing short programs & Prediction of performance on test exercises & Predicts test performance well & Needs hand-authored rule model per domain\\\hline
Work 02 \cite{piech2015} Piech et al., 2015 & Predict student performance over time & LSTM-based Deep Knowledge Tracing without expert skill labels & Simulated data, Khan Academy, ASSISTments & AUC, 5-fold cross-validation & AUC 0.86 vs 0.69 (ASSISTments); 0.85 vs 0.68 (Khan) & Predicts but does not choose help; less interpretable\\\hline
Work 03 \cite{rus2015} Rus et al., 2015 & Tailor instruction to each learner & DeepTutor, a web-based conversational ITS & Large-scale experiments with high-school students & Experimental evidence at scale & Conversational ITS shown to work at scale over the web & Pre-LLM; not designed for programming\\\hline
Work 04 \cite{aleven2016} Aleven et al., 2016 & Improve learners' use of on-demand help & Help Tutor feedback; rule-based model of help seeking; review & Classroom study with tutoring systems & Help-seeking behaviour; domain learning & More deliberate help use, lasting; no gain in domain learning & Self-regulated help use remains an open challenge\\\hline
Work 05 \cite{puri2021} Puri et al., 2021 & Lack of large AI-for-code datasets & CodeNet: over 14M samples, 55 languages, with tests & Online programming submissions & Baseline classification and similarity results & Test cases for 98.5\% of samples & No tutoring or learner-state data\\\hline
Work 06 \cite{liffiton2023} Liffiton et al., 2023 & Scalable help without over-reliance & CodeHelp: LLM with guardrails, no solutions & First-year course, 52 students, 12 weeks & Usage patterns, perceptions, instructor reflection & Valued availability and error help; easy to deploy & No per-learner adaptive model reported\\\hline
Work 07 \cite{kazemitabaar2024} Kazemitabaar et al., 2024 & Helpful support without revealing code & CodeAid: conceptual answers, pseudo-code, fix suggestions & 12 weeks, 700 students, about 8,000 uses & Usage logs, surveys, 22 student and 8 educator interviews & Four design considerations for AI assistants & No continuous learner model or similar-learner retrieval reported\\\hline
Work 08 \cite{yang2024} Yang et al., 2024 & Help beginners learn Python & PyTutor: ChatGPT ITS with structured hints & 11 weeks; 35 users vs 36 non-users & Engagement, completion and success rates & Most effective for weak-foundation students; over-reliance risk & Single LLM; hints include coding solutions\\\hline
Work 09 \cite{ma2024} Ma et al., 2024 & Course-specific AI tutoring & RAGMan: retrieval-augmented, homework-specific tutors & Intro course, 455 students, 5 tutors & Usage, in-scope accuracy, perceptions & About half used; 98\% accurate in scope; 78\% said it helped & Retrieves course content, not learner experience; self-report\\\hline
Work 10 \cite{schmucker2024} Schmucker et al., 2024 & Cheaper authoring of conversational tutors & Ruffle\&Riley: scripts from lesson text; two LLM agents & Biology lessons; 2 online studies, 200 participants & Engagement, perceptions, learning gain & Engaging and helpful; no significant learning difference vs reading & Not programming; engagement is not learning\\\hline
Work 11 \cite{zhan2025} Zhan et al., 2025 & Evaluate tutoring with scarce data & CoderAgent: ACT-R-inspired LLM simulation of programming & CodeNet Python subset; CSEDM Java & Accuracy, CodeBLEU, AUC on next-step tasks & Interpretable, accurate simulation & Simulation, not live tutoring\\\hline
Work 12 \cite{wang2025} Wang et al., 2025 & Proactive guidance in coding tutoring & TRAVER: knowledge tracing plus turn-by-turn verification; DICT protocol & Simulated students; code-generation tests & Task success rate & Significantly higher success rate & Simulated evaluation; no help-amount decision or retrieval\\\hline
Work 13 \cite{pitts2025} Pitts et al., 2025 & Scale personalised novice support & Survey of feedback, assessment and knowledge modelling & Recent LLM programming-education studies & Not applicable (survey) & Human-in-the-loop and scaffolding work best & Not reported (count of works); proposes no system\\\hline
\end{longtable}\endgroup\end{landscape}
